\documentclass[11pt]{article}
\usepackage[margin=1in]{geometry}
\usepackage{amsmath,amssymb,bm}
\title{The battery-only charge hypothesis}
\author{}
\date{}
\begin{document}
\maketitle

\section{Problem Setup}
The presentation asks whether charge at the battery alone can make an
electric field that guides current around a bent circuit. For this
comparison, assume that the two battery terminals carry uniform,
opposite charge densities. Put no charge on the wire surfaces or
resistor interfaces. This is deliberately not the conducting
circuit's solution.

As in the rest of the app, the geometry is constant in the $z$
direction. A terminal edge in the $xy$ plane extends indefinitely
along $z$. We calculate its field at the center of every triangle
of the unified mesh.

\section{Finding the field of a charged terminal edge}
Start with a line charge of strength $\lambda$ per unit $z$ length.
A circular Gaussian surface of radius $r$ and unit depth has
outward electric flux $E(2\pi r)$. Gauss's law therefore gives
$E=\lambda/(2\pi\epsilon r)$. In vector form, the field at
$\bm r$ from charge located at $\bm r'$ is proportional to
$(\bm r-\bm r')/|\bm r-\bm r'|^2$.

For a terminal mesh edge $e$ with constant charge density
$\sigma_e$, a small piece $ds'$ contributes line charge
$\sigma_e\,ds'$. Add all pieces of the edge:
\begin{equation}
 \bm E_e(\bm r)=
 \frac{\sigma_e}{2\pi\epsilon}
 \int_e
 \frac{\bm r-\bm r'}{|\bm r-\bm r'|^2}\,ds'.
 \label{eq:edge}
\end{equation}
The total toy field is the sum of equation~\eqref{eq:edge}
over both terminals. The app evaluates each edge integral
numerically with 12-point Gauss--Legendre quadrature. The
evaluation points are triangle centers, away from the charged
edges. This field spreads into the air and becomes weaker farther
away; it does not respond to the shape of the rest of the wire.

\section{Choosing the amount of terminal charge}
The toy hypothesis does not tell us the charge density. We choose
one scale so that its field inside the battery is comparable with
the actual battery field. Let the battery EMF be $\mathcal E$,
its length $\ell_b$, and its driving direction $\bm d_b$.
The conduction law inside the battery is
\begin{equation}
 \bm J=\frac{\bm E+\bm F}{\rho_b},\qquad
 \bm F=\frac{\mathcal E}{\ell_b}\bm d_b.
\end{equation}
In the limit of very small battery resistivity $\rho_b$ and finite
current, $\bm E$ approaches $-\bm F$. We use the axial value
$-\mathcal E/\ell_b$ as the calibration target.

First give one terminal density $+1$ and the other $-1$, and
calculate the resulting field $\bm E^{(1)}$. For the battery
triangles $B$, compute its area-weighted response along the
battery direction:
\begin{equation}
 R=\frac{\sum_{h\in B}A_h\,
                  \bm E_h^{(1)}\cdot\bm d_b}
         {\sum_{h\in B}A_h}.
\end{equation}
Electric field is linear in charge density, so multiplying both
terminal densities and the trial field by
\begin{equation}
 \alpha=\frac{-\mathcal E/\ell_b}{R}
\end{equation}
makes the battery-average axial field match the target.
The app displays $\bm E_{\rm toy}=\alpha\bm E^{(1)}$ and
the corresponding terminal charges. This calibration does not
make the toy field satisfy the wire's boundary conditions.

\paragraph{Code:} \texttt{calculations/battery\_only.py}.
\end{document}
